\chapter[Sermon 88. What Those Thousand Years Shall Be, and of the Certain Felicity of Souls After the Death Corporal, and of the First Resurrection, and Second Death.]{What Those Thousand Years Shall Be, and of the Certain Felicity of Souls After the Death Corporal, and of the First Resurrection, and Second Death.}
\label{sermon:088}
\markright{Sermon 88. What Those Thousand Years Shall Be, and of the Certain ...}

\index{Resurrection}And I saw thrones, and those who sat on them, and judgment was given to them. And I saw the souls of those who had been beheaded for the testimony of Jesus and for the word of God, who had not worshipped the beast, nor his image, nor received his mark on their foreheads or in their hands. They lived and reigned with Christ for a thousand years. But the rest of the dead did not live again until the thousand years were finished. This is the first resurrection. Blessed and holy is he who has part in the first resurrection. On such, the second death has no power, but they shall be priests of God and of Christ, and shall reign with him for a thousand years.

\index{Antichrist}\index{Millenarianism}By these words John declares himself, explaining what those thousand years will be. Not such, without doubt, as many imagine (among whom are counted the Millenaries or Chiliasts), who say that there will be tranquility on earth, and that during those years the saints here on earth will reign corporally with Christ in exquisite pleasures and joys. For John himself refutes this opinion, while he shows how the saints will be beheaded by the beast and his image, and that those who remain in death will not live again or receive the gospel of Christ. It is manifest, therefore, that the beast and his image will be present during those thousand years. It is evident that the gospel of Christ will shine so brightly during those thousand years that Satan will be so strictly bound that nevertheless not all will receive the gospel, nor will there be quiet tranquility; but that the saints, for Christ’s truth, will suffer persecution by the beast, and that many will not believe the gospel but rather resist it and perish. Yet that the Devil in the meantime will not have so great power as he obtained after the thousand years were finished, nor that the gospel will in those thousand years be so darkened as it was after it was corrupted and depraved. And he touches upon certain opinions that are right notable and necessary, and opens them, namely, what the state will be of those who either are killed for Christ or reject Antichrist; for their souls do not sleep until the judgment, but live with Christ in heaven. He treats moreover of the first resurrection and the second death. Thus, to those who marvel where the souls of the dead will become and what they will do immediately after corporal death, he answers, and declares as much as is necessary to know.

\index{Arethas of Caesarea}\index{Judgment, Last}Therefore John sees thrones, and those who sit on them. And who are those who sit? He adds an explanation and says, “the souls of those who were beheaded.” For it is taken as an explanation, as if to say, those who sit on the heavenly thrones are the souls of those who were beheaded. Souls are not beheaded, but bodies: the souls remain in their state and life. Therefore, he says the souls of those whose bodies were beheaded or slain. And here let us note that John does not speak of the bodies [illegible] at the last judgment, changed, or raised again, but of the souls released from the bodies of the martyrs. For he speaks of souls loosed from the bodies, before the Judgment, according as every one in his time lives here in this world and is called from hence by death. For Arethas, Bishop of Caesarea, also explains this as the souls of martyrs; yet he does not think that no one should be saved unless he die by the tyrants’ sword. For he adds this moreover: or truly he names to be beheaded, tropically, those who have mortified their members on earth. Thus far he. And we also have shown before that first and chiefly the holy martyrs are rewarded with eternal life, secondly all those who have honored God truly and have done penance, and have crucified their flesh with all its concupiscences.

And he states explicitly that the saints were beheaded, not for theft, murder, or wrongdoing, as Saint Peter also teaches, 1 Peter 4. But for the word of God and the testimony of Jesus Christ. The word of God is the very Son of God our Savior, and the testimony is that wholesome gospel, and the very preaching and professing of the same, as we have explained previously through the examination of Scripture. They are also numbered among the saints those who have not worshipped the beast, and so on. And these are the martyrs beheaded or slain, because they worshipped God, but they would not worship the beast or its image. However, not all who reject Antichrist are slain, and therefore he particularly mentioned them as a distinct member. But what it means to worship the beast and its image, and to receive its mark, and so on, I have explained at length in chapter 13.

Now let us see what their condition is, those who shed their blood for Christ and abhor Antichrist with all his enchantments: they live, he says, that is, by faith in this present world. As Paul also said: I no longer live, but Christ lives in me. And from that same life follows everlasting life in another world. Therefore John added, “and they reigned with Christ for a thousand years,” meaning that entire span of time. Not that they did not reign and live with Christ afterward, but that their souls until the judgment have not slept, but rather have lived a blessed life in heaven. Which from the beginning he indicates by another vision. For he sees a seat set, and the souls sitting on them. And by a figurative expression he signifies that certain seats and honorable places are prepared in heaven for the blessed souls, as the Lord himself says in the gospel: “In my Father’s house are many mansions, and I go to prepare a place for you.” He calls the seats thrones, alluding to the royal thrones of kings. But of these heavenly seats, we must conceive of greater, divine, and spiritual matters. They sit not because they do nothing but sit on a cushion, but they reign, triumph, rest, live, and have enjoyment of the comfort, joy, and everlasting glory. This is the manner in which the souls and spirits sit.

He adds moreover that judgment was given to those souls, truly because they are exempted from judgment and do not come into judgment (even as our Savior says), but have passed from death to life. It is also declared in another place, in what sense the saints are said to sit upon the seats and judge the world: where it is evident that all judgment is given to the Son. It is evident therefore by this infallible passage of Scripture that the souls of the saints do not sleep after the death of the body until the last judgment, but live in heaven with Christ. But at the judgment they shall return to their bodies raised again, and together with their bodies shall be received into blessed seats. And this is the condition of the faithful. From this hope let us never allow ourselves to be withdrawn. In my Decades I have discussed more at length about the souls separated from their bodies and have shown that they do not sleep.

\index{Papacy}And here I cannot refrain, but must set forth and recite that which Dr. John Funceius, a learned and diligent man, who has read much in the tenth book of his Chronology, under the year of our Lord 1332, wrote in these words: About this time, the most holy father, Pope John, the twenty-second of that name, fell into this heresy, which he also openly professed and taught, that souls do not see God before the last day. For so his father had taught him, deceived by the visions of Tantalus, which were commonly carried abroad in writing. And Pope John sent two preachers to Paris, namely, a couple of Friars, one of the order of preachers, and another Minorite, to profess his error there. But one Thomas, a preacher of England, resisted the Pope stoutly, whom the Pope committed to prison. And the King of France called a synod in his palace, in the forest Vitinian, where all those assembled subscribed against the Pope. Then the king sent ambassadors to the Pope, exhorting him to recant his error and that he would deliver Thomas out of prison. Which enlarged the prisoner; and also (as it is said), following the admonitions of his friends, at the hour of death repented. So much Funceius. It is a shame therefore for some, who at this day, in so great a light of the gospel, dare renew that most foolish error, affirming that souls separated from their bodies lie snoring, I know not in what dormitory, neither to feel anything, till at the day of judgment they be joined again to their bodies, and rise again.

John adds: and the remnant of the dead did not live again until the thousand years were accomplished. Not that they lived afterward, but that they revived never at all. As Scripture says in another place, Michal, David’s wife, remained barren until the day of her death: not that she had a child after her death. But whom does he mean by the remnant of the dead? Surely all those who descend from Adam are dead. As Paul declares well in chapter 5 to the Romans. But we have heard how some through faith have received Christ, and so being quickened, have shed their blood for Christ, and would not worship the beast or his image. Now is added to this number: but the remnant of the dead, who are neither regenerated through faith nor would bestow their lives for Christ, but would rather worship the beast and his image, these, I say, lived not because of their unbelief. For without faith there is no true life in this world.

We speak nothing here of the vital or natural life. And we say that life is double, or of two sorts: the one spiritual, which is of faith and of the Spirit of God, and of Christ, which is received by faith and lives in the hearts of his people, and his life in them. For the Lord himself says: “He who eats me shall also live for me.” The other life is everlasting, that is, of another world, in which we shall see God as he is, and shall be as he is, living in God and with God forevermore. Conversely, death is of two sorts: spiritual, whereby lacking Christ and his Spirit, and void of faith, we live in sin. The Apostle, speaking of this death, says that a widow living wantonly is dead while alive. And the Lord also to the disciple who wanted to return home and bury his parents, says: “Let the dead bury their dead.” There is also everlasting death, that is, everlasting wretchedness and misery, which follows the spiritual. Yet see what we have said of double death in chapter 3 of this book, in expounding the Epistle to those of Sardis. Therefore, John here signifies that there will be many in these thousand years who will not receive the gospel with a lively faith, and therefore will remain in death, as the Lord said in John 8. Therefore, they err shamefully who suppose that all nations in the whole universal world will come one day to an unity of faith and most assured peace in this life.

\index{John, Saint}\index{Daniel (Book of)}And Saint John himself again explaining his own words says, “This is that first resurrection.” Which I ask you, by which men receive Christ through true faith and rise from sin into a new life. The Apostle speaks much of this in Romans 6. He also speaks, in Ephesians, from Isaiah: “Awake, you who sleep, and rise from the dead, and Christ shall shine upon you.” Therefore, those who do not acknowledge their sins, nor are regenerated, nor are quickened by faith in Christ, nor rise again with Christ in a new life, are not partakers of the first resurrection. The second resurrection is that universal resurrection of all flesh, in which all men will indeed arise, but in different states: the faithful to everlasting life, the unfaithful to everlasting death. The Lord himself also repeated this from Daniel chapter 12, in John chapter 5.

And he demonstrates, by example and in an apostolic manner, a threefold fruit or effect of the first resurrection. First, he says, “Blessed and holy is he who is partaker of the first resurrection.” He is blessed, meaning happy and heir of celestial and eternal life. Holy: that is to say, purified, sanctified, and justified. For faith in Christ does sanctify and make blessed. Therefore, in those who are thus sanctified, the second death has no place nor power. And the first death is the death of sin; therefore, the second death is eternal damnation. See what I have spoken hereof before in chapter 2 of this book, in the Epistle to the church of Smyrna. Finally, the faithful are made priests of God and of Christ, the elect, I mean, segregated, notable, excellent, both of God and Christ most dearly beloved, who in eternal life might offer eternal praises to God. It is repeated again, and they shall reign with him for a thousand years. And this signifies that all saints shall reign with Christ forever, but chiefly the souls, even also before the judgment.

\index{Augustine, Saint}\index{Primasius}Primasius, Bishop of Utica, explaining this passage, says that it is not spoken only of bishops and priests, but just as we call all Christians by reason of the mystical chrism or anointing, so are all priests, for they are members of a Priest, of whom the Apostle, Saint Peter, says, “an holy people, a royal priesthood.” Thus he says. But this whole passage concerning the binding and loosing of the Devil, the thousand years, and the first resurrection and second death, Saint Augustine has well and diligently, for his time and for so much as he could see, discussed at length in Book 20 of \textit{The City of God}. I propose these things of mine to be diligently considered by the faithful. Let every man hold that which he shall think most consonant to the truth. To the Lord our God be praise and glory, now and evermore. Amen.
